% !TeX root = ../papers/nc-p3.tex
\section{階数15の一般水準の全通常点}
\label{sec:rank-fifteen-generic-levels}
\begin{theorem}\label{thm:rank-fifteen-generic-levels}
標数零で、階数15の両表示に共通な有限集合
$E\subset\overline{\mathbb Q}$ が存在し、$\xi\notin E$ の種数5水準では
\emph{全ての通常窓}が標準Koszul左右AS正則である。
各始点は四つの中心四次式の多項式環上で階数256の自由加群である。
特に、通常で非正則な一般窓を持ち、水準を支配する既約曲線はない。
\end{theorem}
\begin{proof}
まず、定理\ref{thm:rank-fifteen-horizontal-specialization}で残った
$c_0^2-c_0+1=0$ の上の $K=\pm1$ の八点を扱う。
$t=n/(c^2-1)$ とする。分母はこれらの点の単元なので、
曲線だけでなく相対曲面でも $n$ と $t$ は同じ主イデアルを生成する。
命題\ref{prop:rank-fifteen-relative-trace-centers}の関係・中心を
$t^{25}$ を法として再構成する。$K(0)=1$ とし、$c$ の係数環では
$c^2-c+1$ と $c^2+c+1$ の\emph{両方}において単元となる分母だけを許す。
係数関数の定数項は $\mathbb Q(c)$ で正確に扱う。

\path{computations/verify_rank_fifteen_generic_levels.py} は級数ごとに
$\mathbb Z[c]$ の一つの共通分母を取り、$t$ の打切り多項式として和と積を取る。
逆元を取るのは定数項の分子が上の二つの既約多項式のどちらでも割れないときである。
単元が尽きて掃き出し中の行を $t$ で割るときは、残る全ての定数項の分子が
$\mathbb Z[c]$ の零多項式であることを調べる。このため任意の通常な一般窓を持つ弧で
同じ行が使える。平方根の初期値、関係格子の構成、長さ5までの未切断の道と
可換条件の核を取る手順は前命題の通りである。
文字0の二式は自由座標 $(1,3,5),(2,7,11)$、文字1の二式は座標基底から取る。
各始点で商に一文字を加え、末尾の二次関係と長さ4での中心式を課す。
各長さで四つの文字ブロックの最大除算回数だけ精度を下げる。
第二分岐の両符号、両表示、全四始点で長さ13まで正の精度を保ち、次元の上界は
\[
 1,4,10,20,31,40,44,40,31,20,10,4,1,0
\]
となる。

$K(0)=-1$ については $(c,K,r,w)\mapsto(-c,-K,-r,-w)$ を使う。
$n$ は同じで、二つの平方根の初期値 $c+K(0),c-K(0)$ はともに反転する。
$z=r^2,\lambda=w/r,(K\pm cn)/r$ は変わらず、第一・第二表示の窓はそれぞれ
まったく同じである。$-c_0$ は $c^2+c+1$ の根なので同時単元判定が適用できる。
二位置移動は第一分岐のもう一方の符号を尽くす。
よって八点に特殊化するどの通常な一般窓にも四中心商の証明書が得られる。
各次数で一般窓から来るHilbert次元の下界にこの上界を合わせると
四中心上の自由性と階数256が従う。通常な切断の完全対合を用いた
定理\ref{thm:klein-rational-period-four}の左右Koszul複体の議論が適用できる。

最後に四中心商の証明書を持つ通常窓の集合は構成可能である。実際、長さ5までの
経路基底と交換子行列の階数で有限層化できる。四式の係数は有限次元の核に属し、
四始点・長さ13までの各上界は有限行列の階数条件である。
よって係数の選択を含む構成可能集合を射影すればよい。
水準を支配する曲線の一般点には、この証明の八点の場合と
定理\ref{thm:rank-fifteen-horizontal-specialization}の証明の他の全ての場合により、
実際にこの有限証明書がある。従って幾何学的一般ファイバーのどの閉点も
証明書を持つ。もしその補集合と一般ファイバーとの交わりが非空なら、
構成可能性によりそこに閉点があるので矛盾する。補集合の $c$ 直線への像も
構成可能で一般点を含まないため有限である。この議論を
$\overline{\mathbb Q}$ 上で行い、$L^2=1+c^2$ の選択を有限射で降ろし、
両表示の例外の $\xi=c^2$ における像の和集合を取ればよい。
基礎体が代数閉でないときも、中心の係数空間における良い四式の非空開集合は
無限な基礎体上に点を持ち、左右の完全性は体拡大から降下する。
\end{proof}
